% ---
% title: Calculate Log Probability of Diffusion
% date: 2025/12/7
% tag: Diffusion
% ---
\definecolor{mdc1}{HTML}{FFFFFF}
\definecolor{mdc2}{HTML}{DD0000}
\section{推导}

流匹配里，策略模型 \(θ\) 下的条件对数概率 \(logp_θ(x_{t−1}∣x_t,c)\) 用高斯概率公式计算：

\[
\log p _ {\theta} \left(x _ {t - 1} \mid x _ {t}, \boldsymbol {c}\right) = - \frac {\left\| x _ {t - 1} - \mu_ {\theta} \left(x _ {t} , t\right) \right\| ^ {2}}{2 \sigma_ {t} ^ {2} d t} - C _ {t}, \tag {6}
\]

其中 \(x_{t - 1} = \mu_{\theta_{old}}(x_t,t) + \sigma_t\sqrt{dt}\cdot \pmb {\epsilon},\pmb {\epsilon}\sim \mathcal{N}(0,\pmb {I})\)，\(C_t\) 是常数。

于是后验重要性比例的对数 \(\log r_t(\theta)\) 可以写成下面这样。为简化，记 \(\Delta \mu_{\theta} = \mu_{\theta_{old}}(x_t,t) - \mu_{\theta}(x_t,t)\)：

\begin{align}
\log r_t(\theta)
&= \log p_{\theta}(x_{t-1} \mid x_t, \boldsymbol{c})
 - \log p_{\theta_{\text{old}}}(x_{t-1} \mid x_t, \boldsymbol{c}) \\
&= - \frac{\| \mu_{\theta_{\text{old}}}(x_t, t) - \mu_{\theta}(x_t, t) + \sigma_t \sqrt{dt}\cdot \boldsymbol{\epsilon} \|^2}{2 \sigma_t^2 dt} \\
&\quad + \frac{\| \mu_{\theta_{\text{old}}}(x_t, t) - \mu_{\theta_{\text{old}}}(x_t, t) + \sigma_t \sqrt{dt}\cdot \boldsymbol{\epsilon} \|^2}{2 \sigma_t^2 dt} \\
&= - \frac{\| \Delta \mu_\theta + \sigma_t \sqrt{dt}\cdot \boldsymbol{\epsilon} \|^2}{2 \sigma_t^2 dt}
 + \frac{\| \boldsymbol{\epsilon} \|^2}{2} \\
&= - \frac{\| \Delta \mu_\theta \|^2}{2 \sigma_t^2 dt}
 - \frac{\Delta \mu_\theta \cdot \boldsymbol{\epsilon}}{\sigma_t \sqrt{dt}}
\tag{7}
\end{align}

\(\epsilon \sim \mathcal{N}(0,I)\)，下面用一维高斯说明（不失一般性）：

\[
\mathbb{E}_{\epsilon \sim \mathcal{N}(0,I)}[\log r_t(\theta)] = -\frac{\|\Delta\mu_\theta\|^2}{2\sigma_t^2\Delta t}
\]

这里和 LLM 有一个\colorbox{mdc2}{\textcolor{mdc1}{关键区别}}：语言模型算的是离散 token 的概率，扩散模型算的是高斯状态转移的概率。\colorbox{mdc2}{\textcolor{mdc1}{二次项给对数重要性比带来一个随时间步变化的负偏置}}，如图 2(b)。期望比率通常小于 1，正优势的样本很少能超过上裁剪边界。于是来自过度自信的正预测的梯度大多被保留，负样本的梯度反而被约束得更紧，策略容易过优化。

重要性比例的方差还取决于去噪调度器参数 \(\sigma_t\) 和 \(dt\)，在不同时间步上差别很大。低噪声步骤里裁剪阈值经常被超过，高噪声步骤里几乎不触发，裁剪的不平衡被放大，策略会朝时间步特定的方向过优化。

[\textasciicircum{} 来源]:FlowGRPO
